% Page 068 — page imprimée 64 bis
% La vallée de la Jonte

\begin{center}
\fbox{\parbox{0.65\textwidth}{
\textbf{Meyrueis se trouve au confluent de trois vallées,\\
Jonte, Brèze et Bétuzon. Sans doute faut-il chercher\\
dans cette situation l’étymologie de son nom\\
« au milieu des rivières ».
}}}

\vspace{1em}

\textbf{Je vous propose de les remonter, l’une après l’autre.}

\vspace{1.5em}

{\Large\bfseries LA VALLEE DE LA JONTE}
\end{center}

\medskip

C’est au pied du signal des Fons, au bas de la pelouse de l’Aigoual, près de la grande congère
qu’elle prend sa source. Ruisselet ruisseau, rivière elle cascade dans les gorges étroites des
Escarabits avant de baigner les maisons de Jontanels et de s’épanouir du côté de Gatuzières
abandonnant le granite et les schistes de sa naissance pour serpenter sous les hautes couronnes
calcaires du Méjean.

\medskip

Après Meyrueis elle sépare désormais Causse Noir et Causse Méjean ; les gorges deviennent
profondes abritant bien des grottes, celles de la Chèvre, de la Vigne et de Nabrigas, arcade des
bergers et, si bien cachée, la grotte Amélineau aux mille stalactites fistuleuses éblouissantes,
la perte des Hérans et la résurgence des Douze, le Roc Saint Gervais, et l’ermitage Saint
Michel avant de s’abandonner dans les eaux vertes du Tarn pour gagner l’Océan.

\medskip

Maintenant tel un pêcheur de truites, au grand dam des géographes, nous remontons son cours.
Barrée de nombreuses chaussées, la Jonte arrosait prés et jardins et faisait tourner de
nombreux moulins. L’un des plus anciens, le moulin de Capelan fut équipé postérieurement
d’une petite turbine qui permit l’électrification du village dans les années 1900, le moulin de
Malbois rue des Apiès, au centre du village, un des derniers à faire de la farine avec celui de
Gatuzières jusqu’en 1950-1960 environ. (22 moulins sur nos trois rivières vers 1850 !!!).

\medskip

Sur la rive gauche de la Jonte, à notre droite donc, après « l’arbre du poivre » (voir dans « La vallée
de la Brèze) et la piscine se découpent les maisons d’Ayres et invisible, sous de grands arbres
entourés d’un haut mur se cache le Château d’Ayres dont nous connaissons l’histoire récente.

\medskip

Son moulin, le Moulin d’Ayres que l’on rejoint, à pied, par le Pont de Six Liards a connu bien
des évolutions, résumées dans les pages suivantes.

\medskip

Nous le quitterons pour remonter au plus près le cours de cet affluent du Tarn.

\medskip

Enfin tout près de la source de la Jonte, l’histoire de Cabrillac a été relatée par Yves Grellier
en 1991 « Cabrillac : Un hameau sur L’Aigoual ».
